\documentclass[11pt]{article}
% =====================================================================
%  Status:      research note (exploratory).  NOT a paper claim.
%  Provenance:  derived in the session of 2026-09-28/29; the companion
%               note 07 fixes the group G and its coordinate conventions.
%  Promotion:   nothing here is promoted.  See §7.
% =====================================================================
\usepackage[utf8]{inputenc}
\usepackage[T1]{fontenc}
\usepackage{amsmath}
\usepackage{amssymb}
\usepackage{amsthm}
\usepackage{booktabs}
\usepackage{hyperref}
\hypersetup{colorlinks=true,linkcolor=blue,urlcolor=cyan}
\usepackage[a4paper,margin=1in]{geometry}

\newtheorem{theorem}{Theorem}[section]
\newtheorem{lemma}[theorem]{Lemma}
\newtheorem{proposition}[theorem]{Proposition}
\newtheorem{corollary}[theorem]{Corollary}
\newtheorem{definition}[theorem]{Definition}
\newtheorem{remark}[theorem]{Remark}
\newtheorem{example}[theorem]{Example}

% Long paths and inline formulas in \texttt cannot hyphenate.
\tolerance=1000
\emergencystretch=2.5em

\title{Isometric realisations of the joint group in the eight Thurston geometries}
\author{Mingli Yuan and DeepSeek Harness Agent}
\date{2026-09-29}

\begin{document}
\maketitle

\begin{abstract}
The joint group $G=\mathbb R\times\operatorname{Aff}^+(1,\mathbb R)$ of the
companion note carries a three-parameter family of left-invariant metrics. This
note asks which of Thurston's eight three-dimensional model geometries admit a
\emph{lossless} realisation of $G$, meaning a smooth homomorphism
$G\to\operatorname{Isom}(X)$ whose orbit map is a global diffeomorphism. Exactly
two do: $\mathbb H^2\times\mathbb R$ and $\widetilde{\mathrm{SL}}_2(\mathbb R)$.
For the other six, every smooth isometric action kills either the derived
subgroup or the centre. The second case is the interesting one for the programme:
after a Reeb unit extension and an inversion, the contact model of Paper 0 §9 is
exactly the $\widetilde{\mathrm{SL}}_2(\mathbb R)$ geometry, with the horizontal
metric $h$ of Paper 0 as its contact distribution metric. \S\ref{sec:boundaries}
records what is not claimed, in particular that $G$ is \emph{not} isomorphic to
$\widetilde{\mathrm{SL}}_2(\mathbb R)$.
\end{abstract}

\tableofcontents

\section{Status and provenance}
\label{sec:status}

\textbf{Status.} Research note, exploratory. The classification theorem below is
proved here, but $G$ is not a Paper 0 object and nothing here is a paper claim or
a native certificate. The identification in \S\ref{sec:paper0} is the item most
likely to matter to the programme; it is a correspondence of structures, and it
is stated with its conventions explicit.

\textbf{Provenance.} Written 2026-09-29. Depends on the companion note
\path{notes/foundations-and-geometry/07-joint-quotient-acs-charge-aes-operator.tex}
for the group law, the centre and derived subgroup, the projections, and the
contact form. Sources read in \path{paper-0/}: \path{sections/09-contact-of-tearing.tex}
(the contact structure, the horizontal fields, the curvature, the scalar
horizontal differential).

\textbf{Not read for this note.} The classification of the eight model
geometries and their isometry groups is cited as standard; it is not re-derived
here. Only the obstructions \emph{for this particular $G$} are proved below.

\section{The question}
\label{sec:question}

Throughout,
\[
G=\mathbb R\times\operatorname{Aff}^+(1,\mathbb R),\qquad
g_2g_1=(A_2+A_1,\ M_2+M_1,\ B_2+e^{M_2}B_1),
\]
with inverse $g^{-1}=(-A,-M,-e^{-M}B)$, and
\[
\mathfrak g=\operatorname{span}\{Z,H,N\},\qquad
[Z,H]=[Z,N]=0,\qquad [H,N]=N,
\]
where $Z$ is the $A$-direction, $H$ the $M$-direction and $N$ the $B$-direction.
By the companion note,
\[
[\mathfrak g,\mathfrak g]=\mathbb RN,\qquad \mathfrak z(\mathfrak g)=\mathbb RZ,
\]
so both the derived subalgebra and the centre are one-dimensional and
characteristic.

\begin{definition}[Lossless geometric realisation]
\label{def:lossless}
Let $X$ be a Thurston model geometry. A \emph{lossless realisation} of $G$ in
$X$ is a smooth homomorphism $\rho:G\to\operatorname{Isom}(X)$ together with a
basepoint $o\in X$ such that the orbit map
\[
F:G\longrightarrow X,\qquad F(g)=\rho(g)o
\]
is a global diffeomorphism.
\end{definition}

Since $\dim G=3=\dim X$, $F$ a diffeomorphism is equivalent to the action being
free and transitive. Note also that
\[
F(g_2g_1)=\rho(g_2)F(g_1),
\]
so a lossless realisation lets every point of $X$ be read as a unique group
element, with the group law corresponding to the action.

\begin{remark}[Why ``lossless'' and not merely ``faithful'']
\label{rem:why}
A faithful action with a large isotropy is not lossless: one point no longer
recovers a group element. Definition~\ref{def:lossless} requires both freeness
and transitivity, i.e.\ that the orbit map be injective \emph{and} surjective.
\end{remark}

\section{The classification}
\label{sec:classification}

\begin{theorem}
\label{thm:classification}
A lossless realisation of $G$ exists in exactly two of the eight model
geometries:
\[
\mathbb H^2\times\mathbb R
\qquad\text{and}\qquad
\widetilde{\mathrm{SL}}_2(\mathbb R).
\]
In each of the other six, no smooth homomorphism $G\to\operatorname{Isom}(X)$ is
injective.
\end{theorem}

\begin{proof}[Proof of the positive half]
See \S\ref{sec:two}. The formulas are written once and realise both geometries,
which also proves that the action is free and transitive in the two cases.
\end{proof}

The negative half is proved case by case. The two lemmas below do most of the
work.

\begin{lemma}[No real eigenvalue $1$]
\label{lem:no-eigenvalue-one}
Let $\mathfrak i$ be the isometry Lie algebra of $\mathbb E^3$, $\mathbb S^3$,
$\mathbb S^2\times\mathbb R$ or $\mathrm{Nil}$. Then for every $X\in\mathfrak i$,
every complex eigenvalue of $\operatorname{ad}_X$ is either $0$ or purely
imaginary.
\end{lemma}

\begin{proof}
$\mathfrak{so}(3)$ and $\mathfrak{so}(4)$ are compact, so $\operatorname{ad}_X$
is skew-adjoint with respect to an invariant inner product and has purely
imaginary spectrum; adding a central direction adds only zero eigenvalues. For
$\mathfrak{isom}(\mathbb E^3)=\mathfrak{so}(3)\ltimes\mathbb R^3$, write
$X=(\omega,v)$ and $(\eta,u)$ for a second element; then
$\operatorname{ad}_X(\eta,u)=(\omega\times\eta,\ \omega\times u-\eta\times v)$,
which is block upper triangular with diagonal blocks $\omega\times$, so the
spectrum is that of $\omega\times$ repeated twice, namely $0,\pm i\|\omega\|$.
For $\mathfrak{isom}(\mathrm{Nil})=\mathfrak{heis}_3\rtimes\mathfrak{so}(2)$ use
the flag
\[
0\subset\mathfrak z(\mathfrak{heis}_3)\subset\mathfrak{heis}_3
\subset\mathfrak{isom}(\mathrm{Nil});
\]
the adjoint action on the three successive quotients is zero, a planar rotation,
and zero, giving spectrum $0,0,\pm ir$.
\end{proof}

\begin{proposition}
\label{prop:four-killed}
If $X$ is any of $\mathbb E^3$, $\mathbb S^3$, $\mathbb S^2\times\mathbb R$,
$\mathrm{Nil}$, then every smooth homomorphism $\rho:G\to\operatorname{Isom}(X)$
kills $N$, hence the whole one-parameter subgroup $\{(0,0,B)\}$.
\end{proposition}

\begin{proof}
Let $\phi=d\rho$ and write $X_0=\phi(H)$, $Y_0=\phi(N)$. Then
\[
[X_0,Y_0]=\phi([H,N])=\phi(N)=Y_0 .
\]
If $Y_0\ne0$ then $Y_0$ is an eigenvector of $\operatorname{ad}_{X_0}$ with real
eigenvalue $1$, contradicting Lemma~\ref{lem:no-eigenvalue-one}. Hence
$Y_0=0$, and since $G$ is connected the corresponding one-parameter subgroup is
killed.
\end{proof}

\begin{proposition}
\label{prop:h3-killed}
Every smooth homomorphism $\rho:G\to\operatorname{Isom}(\mathbb H^3)$ kills the
centre $Z$; hence none is injective.
\end{proposition}

\begin{proof}
$\mathfrak{isom}(\mathbb H^3)\cong\mathfrak{sl}_2(\mathbb C)$ as a real Lie
algebra. Put $X_0=\phi(H)$, $Y_0=\phi(N)$, so $[X_0,Y_0]=Y_0$. If $Y_0=0$ we are
done. Otherwise $X_0$ is not nilpotent (a nilpotent cannot have $1$ as an
eigenvalue of its adjoint action), so its two eigenvalues $\pm\lambda$ are
distinct and nonzero; conjugating, take $X_0=\operatorname{diag}(\tfrac12,-\tfrac12)$.
Then $\operatorname{ad}_{X_0}$ acts on $\mathfrak{sl}_2(\mathbb C)$ with
eigenvalues $0,\pm1$, and $Y_0$ lies in the eigenvalue-$1$ eigenspace, so
$Y_0=cE_{12}$ with $c\ne0$. Now $Z$ is central in $\mathfrak g$, so $\phi(Z)$
commutes with both $X_0$ and $Y_0$. Commuting with $X_0$ forces $\phi(Z)$ to be
diagonal; commuting with $cE_{12}$ then forces its two diagonal entries to be
equal; being traceless, it is $0$.
\end{proof}

\begin{proposition}
\label{prop:sol-killed}
No smooth homomorphism $G\to\operatorname{Isom}(\mathrm{Sol})$ is injective.
\end{proposition}

\begin{proof}
$\operatorname{Isom}_0(\mathrm{Sol})$ is the Sol left-translation group, with
\[
[T,U]=U,\qquad [T,V]=-V,\qquad [U,V]=0 .
\]
Its derived subalgebra is $\operatorname{span}\{U,V\}$, of dimension $2$, whereas
$\dim[\mathfrak g,\mathfrak g]=1$. An injective $\phi:\mathfrak g\to\mathfrak{sol}$
is an isomorphism because both are three-dimensional, and an isomorphism carries
the derived subalgebra onto the derived subalgebra; $1=2$ is a contradiction.
\end{proof}

\begin{remark}
The two groups are not to be identified merely because both involve exponential
scaling: the obstruction is the dimension of the derived subalgebra.
\end{remark}

\begin{proof}[Completion of Theorem~\ref{thm:classification}]
The six negative cases are Propositions~\ref{prop:four-killed},
\ref{prop:h3-killed} and \ref{prop:sol-killed}: in each, no injective smooth
homomorphism exists, so \emph{a fortiori} no lossless realisation does. The
remaining two models are realised in \S\ref{sec:two}.
\end{proof}

\begin{remark}[Reality of the two realisations]
In both positive cases the realisation is not merely a subgroup isomorphism: the
orbit map is a diffeomorphism, so the geometry provides a global coordinate
system in which the group law is the action. This is what makes the phrase
``lossless'' meaningful, and it is also why the classification is stronger than
asking for a faithful action.
\end{remark}

\section{The two realisations share one formula}
\label{sec:two}

Set
\[
t=A,\qquad x=B,\qquad y=e^M>0 .
\]
If the finite verification chain stores $S=e^M>0$ rather than $M$, no
exponential or logarithm has to be evaluated: the map and its inverse are just
$(A,S,B)\leftrightarrow(t,x,y)=(A,B,S)$.

\begin{proposition}
\label{prop:realisation}
Define $F(A,M,B)=(t,x,y)=(A,B,e^M)$ and let $g=(a,m,b)$ act by
\[
\rho(g)(t,x,y)=\bigl(t+a,\ e^mx+b,\ e^my\bigr).
\]
Then $\rho(g_2)\rho(g_1)=\rho(g_2g_1)$, and with $o=(0,0,1)$ the orbit map
$F(g)=\rho(g)o$ is a diffeomorphism. The following three one-forms are
$\rho$-invariant:
\[
\eta_0=dt,\qquad \eta_1=\frac{dy}{y},\qquad \eta_2=\frac{dx}{y},
\]
with structure equations
\begin{equation}
\label{eq:structure}
d\eta_0=d\eta_1=0,\qquad d\eta_2=-\eta_1\wedge\eta_2 .
\end{equation}
\end{proposition}

\begin{proof}
Substituting $(t+a,\,e^mx+b,\,e^my)$ into the three forms gives invariance, since
$d(e^mx+b)=e^mdx$ and $e^my$ scales the denominator equally. The composition law
is term-by-term substitution. The orbit map is $(a,m,b)\mapsto(a,b,e^m)$, a
diffeomorphism onto $\mathbb R^2\times\mathbb R_{>0}$. For the structure
equations, the dual frame is $e_0=\partial_t$, $e_1=y\partial_y$, $e_2=y\partial_x$,
with $[e_1,e_2]=e_2$ and all other brackets zero; the identity
$d\eta^k=-\sum_{i<j}c^k_{ij}\eta^i\wedge\eta^j$ gives \eqref{eq:structure}.
\end{proof}

\begin{remark}
Equation~\eqref{eq:structure} is the abstract bracket $[H,N]=N$ written in the
coframe dual to $\{\eta_0,\eta_1,\eta_2\}$: $\eta_0\leftrightarrow Z$,
$\eta_1\leftrightarrow H$, $\eta_2\leftrightarrow N$.
\end{remark}

Two left-invariant metrics on the same space are now available:
\begin{align}
g_{\mathrm{prod}}&=\frac{dx^2+dy^2}{y^2}+dt^2,
\label{eq:prod}\\[2pt]
g_{\mathrm{SL}}&=\frac{dx^2+dy^2}{y^2}+\Bigl(dt-\frac{dx}{y}\Bigr)^2 .
\label{eq:sl}
\end{align}

\begin{proposition}
\label{prop:two-metrics}
$\rho$ is an isometry for both \eqref{eq:prod} and \eqref{eq:sl}. Pulling back
along $F$,
\[
F^*g_{\mathrm{prod}}=dA^2+dM^2+e^{-2M}dB^2,
\qquad
F^*g_{\mathrm{SL}}=dM^2+e^{-2M}dB^2+\bigl(dA-e^{-M}dB\bigr)^2 .
\]
The first is the product metric of $\mathbb H^2\times\mathbb R$; the second is
the standard $\widetilde{\mathrm{SL}}_2(\mathbb R)$ geometry, i.e.\ the
universal cover of the Sasaki metric on the unit tangent bundle of $\mathbb H^2$,
with the fibre angle taken in $\mathbb R$ and not modulo $2\pi$.
\end{proposition}

\begin{proof}
Invariance of \eqref{eq:prod} and \eqref{eq:sl} follows from
Proposition~\ref{prop:realisation}, since both are built from
$\eta_1^2+\eta_2^2$, $\eta_0^2$ and $\eta_0\eta_2$. The pullbacks are the
substitutions $dt=dA$, $dx=dB$, $dy=e^MdM$.
\end{proof}

\begin{proposition}[The same group does not determine the geometry]
\label{prop:tau-family}
For $\tau\in\mathbb R$ put
\[
g_\tau=dM^2+e^{-2M}dB^2+\bigl(dA-\tau e^{-M}dB\bigr)^2 .
\]
Every $g_\tau$ is left-invariant, and with the orthonormal frame
$E_1=H$, $E_2=N+\tau Z$, $E_3=Z$ one has
\[
[E_1,E_2]=E_2-\tau E_3,\qquad [E_1,E_3]=[E_2,E_3]=0,
\]
\[
K(E_1,E_2)=-1-\tfrac34\tau^2,\qquad
K(E_1,E_3)=K(E_2,E_3)=\tfrac14\tau^2,
\]
and Ricci eigenvalues $\bigl(-1-\tfrac12\tau^2,\ -1-\tfrac12\tau^2,\ \tfrac12\tau^2\bigr)$.
The values $\tau=0$ and $\tau=1$ recover \eqref{eq:prod} and \eqref{eq:sl}.
\end{proposition}

\begin{proof}
Left-invariance is immediate from Proposition~\ref{prop:realisation}. The bracket
$[E_1,E_2]=[H,N+\tau Z]=N=E_2-\tau E_3$ uses $[H,N]=N$ and $[H,Z]=0$. The
sectional curvatures and the Ricci tensor follow from the Koszul formula in an
orthonormal left-invariant frame with these constant structure constants.
\end{proof}

So the product model has no positive sectional curvature, while the contact
model has positive vertical sectional curvature: the two are genuinely different
Riemannian geometries on one and the same Lie group. The extra data that selects
between them is the metric, the contact distribution, the connection and the unit
convention.

\section{The identification with Paper 0's contact model}
\label{sec:paper0}

This is the part of the note that touches the programme directly. Paper 0 §9
works with
\[
\alpha=dB-dA-B\,dM,\qquad \mathcal D=\ker\alpha,\qquad
h=\bigl(dA^2+dM^2\bigr)\big|_{\mathcal D},
\]
and the AES projection $Q(A,M,B)=(-Be^{-M},e^{-M})$.

\begin{lemma}
\label{lem:paper0-data}
Write $\beta=dB-B\,dM$. Then
\[
Q^*g_{\mathbb H^2}=dM^2+\beta^2,
\qquad
Q^*g_{\mathbb H^2}\big|_{\mathcal D}=h .
\]
The Reeb field of $\alpha$ is $R_\alpha=-\partial_A$, and the Riemannian
extension that makes the Reeb direction unit and orthogonal to $\mathcal D$ is
\[
g_{\mathrm{contact}}:=Q^*g_{\mathbb H^2}+\alpha^2
=dM^2+\bigl(dB-B\,dM\bigr)^2+\bigl(dB-dA-B\,dM\bigr)^2 .
\]
It satisfies $g_{\mathrm{contact}}(R_\alpha,R_\alpha)=1$ and
$g_{\mathrm{contact}}(R_\alpha,\mathcal D)=0$.
\end{lemma}

\begin{proof}
$Q$ has components $x=-Be^{-M}$, $y=e^{-M}$, so $dx=-e^{-M}\beta$ and
$dy=-e^{-M}dM$; hence $(dx^2+dy^2)/y^2=\beta^2+dM^2$. On $\mathcal D$ we have
$\alpha=0$, i.e.\ $dB=dA+B\,dM$, so $\beta=dA$ and the restriction is $h$. For
the Reeb field, $\alpha(c\partial_A)=-c$, so $c=-1$; and
$d\alpha=dM\wedge dB$ annihilates $\partial_A$. Orthogonality and unit length are
then immediate from $\alpha|_{\mathcal D}=0$ and $\beta(\partial_A)=0$.
\end{proof}

\begin{theorem}[The contact model is the $\widetilde{\mathrm{SL}}_2$ geometry]
\label{thm:identification}
Let $I(A,M,B)=(-A,-M,-Be^{-M})$ be the inversion. Then
\[
I^*\alpha=dA-e^{-M}dB=:\omega,
\qquad
I^*g_{\mathrm{contact}}
=dM^2+e^{-2M}dB^2+\omega^2
=F^*g_{\mathrm{SL}} .
\]
Consequently Paper 0's contact model, together with the horizontal metric $h$,
after the Reeb unit extension and the inversion, \emph{is} the
$\widetilde{\mathrm{SL}}_2(\mathbb R)$ geometry, with $\mathcal D$ as its contact
distribution.
\end{theorem}

\begin{proof}
With $(A',M',B')=(-A,-M,-Be^{-M})$ one has $dA'=-dA$, $dM'=-dM$ and
$dB'=-e^{-M}dB+Be^{-M}dM$. Hence
\[
I^*\alpha=dB'-dA'-B'dM'
=-e^{-M}dB+Be^{-M}dM+dA-Be^{-M}dM
=dA-e^{-M}dB,
\]
and
\[
I^*\beta=dB'-B'dM'=-e^{-M}dB+Be^{-M}dM-Be^{-M}dM=-e^{-M}dB .
\]
Substituting into $g_{\mathrm{contact}}=dM^2+\beta^2+\alpha^2$ gives
$dM^2+e^{-2M}dB^2+\omega^2$, which is $F^*g_{\mathrm{SL}}$ by
Proposition~\ref{prop:two-metrics}.
\end{proof}

\begin{remark}[The convention switch is not optional]
\label{rem:convention}
Inversion exchanges left and right. In the original coordinates
$\alpha,h,g_{\mathrm{contact}}$ are preserved by \emph{right} translations
(the companion note's statement \emph{Right translations preserve the data}),
whereas the standard left action on
$\widetilde{\mathrm{SL}}_2(\mathbb R)$ preserves $\omega$. Therefore, as a left
group action, the original model must be written $g\cdot p=pg^{-1}$. Without
this conversion it is not correct to call the left-appended arithmetic step an
isometric symmetry of $\alpha$.
\end{remark}

\begin{proposition}[The new invariants]
\label{prop:invariants}
In the left-action coordinates,
\[
\omega=dA-e^{-M}dB,\qquad d\omega=e^{-M}dM\wedge dB,\qquad
\omega\wedge d\omega=e^{-M}dA\wedge dM\wedge dB\ne0 .
\]
With the base oriented by $dx\wedge dy/y^2$ one has
$d\omega=-dx\wedge dy/y^2$, and along a horizontal lift ($\omega=0$), $dt=dx/y$.
For an oriented closed loop $\gamma=\partial\Sigma$ on the base,
\[
\Delta t=\oint_\gamma\frac{dx}{y}=\int_\Sigma\frac{dx\wedge dy}{y^2},
\]
the signed hyperbolic area. This area does not determine the loop, and \emph{a
fortiori} does not recover the word history.
\end{proposition}

\begin{proof}
Direct computation; the last identity is Stokes' theorem.
\end{proof}

\section{\texorpdfstring{$\mathrm{Nil}$}{Nil} as the graded tangent model}
\label{sec:nil}

Although $G$ has no faithful isometric action on standard $\mathrm{Nil}$
(Proposition~\ref{prop:four-killed}), the Heisenberg algebra appears naturally as
the graded tangent of the contact distribution.

\begin{proposition}
\label{prop:nil-tangent}
In the original contact coordinates put
\[
X=\partial_A+\partial_B,\qquad Y=\partial_M+B\partial_B,\qquad W=[X,Y]=\partial_B .
\]
Then $[Y,W]=-W$ and $[X,W]=0$. Giving $X,Y$ weight $1$ and $W$ weight $2$, the
associated graded Lie algebra is the Heisenberg algebra:
\[
[\bar X,\bar Y]=\bar W,\qquad \bar W\ \text{central}.
\]
\end{proposition}

\begin{proof}
Direct bracket computation. The filtration is the contact filtration of
$\mathcal D$.
\end{proof}

\begin{remark}
The grading discards $[Y,W]=-W$ and the hyperbolic curvature. Moreover the
metric carried by the blow-up limit is the Carnot--Carath\'eodory metric, not the
standard Riemannian metric on $\mathrm{Nil}$; the two do share the Heisenberg
graded group. This is entirely compatible with
Proposition~\ref{prop:four-killed}: the graded tangent is not a realisation of
$G$.
\end{remark}

\section{Boundaries: what is not claimed}
\label{sec:boundaries}

\begin{enumerate}
\item \textbf{$G\not\cong\widetilde{\mathrm{SL}}_2(\mathbb R)$.} The left-hand
side is solvable; the Lie algebra of the right-hand side is simple. The correct
statement is that $G$ is a free transitive \emph{isometry subgroup} of that
three-dimensional Riemannian space, acting by the base positive-affine isometries
together with the commuting fibre translations. The standard $\mathrm{SL}_2$
group law on a point must not be substituted for the group law of $G$; the marked
action $\rho$ has to be kept.
\item \textbf{``Lossless'' is about group elements only.} Theorem~\ref{thm:identification}
transports the invariance and the covariant derivative by an explicit
diffeomorphism whose inverse is written down. It does \emph{not} follow that
every isometry of the target space preserves the extra projections, the
assignment, the vocabulary labels or the verification certificates.
\item \textbf{Coordinate encoding is a much weaker question.} The six
non-spherical models $\mathbb E^3,\mathbb H^3,\mathbb H^2\times\mathbb R,
\widetilde{\mathrm{SL}}_2,\mathrm{Nil},\mathrm{Sol}$ are all diffeomorphic to
$\mathbb R^3$, so any one of them admits a global coordinate system encoding three
real parameters losslessly. Such an encoding does not make the group law an
isometric action, and it is not what Definition~\ref{def:lossless} asks for.
\item \textbf{The classification is about smooth actions.} Discrete subgroups,
finite-word syntax, mixed-chirality marked histories, and complex-parameter
groups are not identified with this Lie group anywhere above.
\item \textbf{No metric is canonical.} Proposition~\ref{prop:tau-family} shows
that the same group carries at least a one-parameter family of genuinely
different left-invariant geometries. Neither $g_{\mathrm{prod}}$ nor
$g_{\mathrm{SL}}$ is forced by the group law, and the companion note's
$dA^2+dM^2+e^{-2M}dB^2$ is a third choice.
\item \textbf{Fibre angles are not reduced.} In the
$\widetilde{\mathrm{SL}}_2(\mathbb R)$ model the fibre angle is a real number, not
an angle modulo $2\pi$. Passing to the unit tangent bundle
$T^1\mathbb H^2\simeq\mathrm{PSL}_2(\mathbb R)$ would forget the winding number;
that step is not taken here.
\end{enumerate}

\begin{thebibliography}{9}
\bibitem{Paper0}
M.~Yuan and DeepSeek Harness Agent,
\emph{Arithmetic Expression Geometry 0: Arithmetic Expression Spaces},
\path{paper-0/}, commit \texttt{d89ceb8f8090482317b750407f623c8dcfb37ea2}.
\bibitem{Note07}
Companion note 07 in this directory,
\emph{The joint quotient of the ACS charge and the AES affine operator}.
\end{thebibliography}

\end{document}
